\section{Efficiency Proofs and Numerical Scope}
For Eqs.~\ref{eq:guaranteed} and \ref{eq:possible}, if $L_d\prec U_c$,
placing $d$ at its lower bound and $c$ at its upper bound makes $c$
dominated. Conversely, a realized dominator implies $L_d\prec U_c$.
If $U_d\prec L_c$, $d$ dominates $c$ throughout the box. Otherwise, place
$c$ at $L_c$ and every rival at its upper bound; no rival dominates $c$.
Independent box coordinates permit these assignments. Strict dominance
retains exact ties.

For Eq.~\ref{eq:prices}, minimizing a linear function on a constrained
rectangle requires only its vertices: admissible box corners and the
intersections with the minimum-gap boundary. Uncoupled coordinates choose
the lower or upper endpoint according to the coefficient's sign. Cost,
burden, and outcome minima can be realized simultaneously for a fixed
rival. The existence of any such rival disproves universal retention;
its absence proves retention. Price calculations use integer upgrade
coefficients and round cancellation residues to 12 decimal places.
Independent linear-programming checks agree to $10^{-10}$.

The population screen follows from Hoeffding's inequality for independent
bounded observations, including nonidentically distributed observations in
fixed strata. A union bound over $K$ comparisons gives simultaneous coverage
of at least $1-\alpha$. If a lower bound is strictly positive for every
rival in some coordinate, none can dominate the candidate. No tail-mean
concentration claim is required. An equal screened vector cannot be
discarded, since an omitted tail coordinate could break the tie.

\begin{table}[t]
\centering\small
\caption{All price radii. Entries give counts in R/I/H profile order.
$B_{95}$ counts original certificates meeting the resampling-frequency
criterion. The positive-gap extension was recorded after the first run.}
\begin{tabular}{@{}lrrrr@{}}\toprule
& \multicolumn{2}{c}{Monotone} & \multicolumn{2}{c}{Positive gaps}\\
Radius \% & $G$ & $B_{95}$ & $G$ & $B_{95}$\\\midrule
\tableinput table_joint_all_rows.tex
\bottomrule\end{tabular}
\end{table}

\section{Engineering and Decision Diagnostics}
\begin{table}[t]
\centering\small
\caption{All planned engineering settings, with 128 finite sampled
candidates per problem. Radius scales the sampled objective range.}
\label{tab:transfer}
\begin{tabular}{@{}lrrrr@{}}\toprule
Problem & Radius & $|G|$ & $|F_0|$ & $|P|$\\\midrule
\tableinput table_transfer_rows.tex
\bottomrule\end{tabular}
\end{table}
RE21 models a four-bar truss; RE22 models a reinforced-concrete beam
and retains its constraint-violation objective \cite{tanabe2020}.
Independent intervals perturb all objectives by $\pm1\%,\pm5\%,\pm10\%$
of their finite sampled ranges, with nonnegative lower bounds.
All \BenchmarkDraws{} perturbation frontiers contain $G$ and lie within
$P$. For the first eight candidates per problem, \BenchmarkCorners{}
total corner matrices with only the second objective uncertain reproduce
the exact frontier intersection and union. \BenchmarkEmptySettings{} of
six settings have empty guaranteed sets. These checks concern finite
tables and chosen bounds, not the continuous problems' global frontiers.

\begin{table}[t]
\centering\small
\caption{Illustrative mean-loss selection within each declared point budget.
Free and selected columns give mean weighted service-hours lost; the
last column gives the selected candidate's $k=4$ frequency.}
\label{tab:budget}
\begin{tabular}{@{}lrrrr@{}}\toprule
Profile & Budget & Free & Selected & Event \%\\\midrule
\tableinput table_budget_rows.tex
\bottomrule\end{tabular}
\end{table}
The budget example selects the lowest observed mean loss within each
profile's point budget. The artifact identifies the portfolios and their
burdens. Another objective or a burden constraint can favor another choice;
zero observed events do not imply zero risk.

\begin{table}[t]
\centering\small
\caption{Frontier sizes over every objective subset of a given dimension,
shown as minimum/median/maximum. Tied optima are retained.}
\begin{tabular}{@{}crrr@{}}\toprule
Dimensions & R & I & H\\\midrule
\tableinput table_ablation_rows.tex
\bottomrule\end{tabular}
\end{table}
Tied minima explain the large intermediate single-objective maximum.
Adding an objective can break ties and remove a previously efficient
candidate, so cardinality need not increase with dimension.

\section{Interpretation and Reproduction}


Each registered assertion specifies a source, population, expected groups,
statistic, unit, quantifier, relation, and bound. Missing groups or unsupported
operations stop verification. The eight correct assertions, three historical
contradictions, and eight single-field mutations were constructed after
reviewing the manuscript. They test specified behavior rather than estimate
general detection accuracy.

For paired mean-loss comparisons, let $D_{hi}$ be the candidate difference
in entry stratum $h$, with weight $w_h=n_h/n$. Then
\begin{equation}
 \widehat{\mathrm{Var}}(\widehat\Delta)
 =\sum_h w_h^2s^2_{D,h}/n_h,\quad
 s^2_{D,h}=s^2_{a,h}+s^2_{b,h}-2s_{ab,h}. \label{eq:pairing}
\end{equation}
Setting $v_h=w_h^2s^2_{D,h}/n_h$ gives approximate degrees of freedom
$(\sum_hv_h)^2/\sum_h[v_h^2/(n_h-1)]$ for the family-adjusted paired
t intervals.

Original raw banks and protocol digests remain unchanged. Added analyses
start from committed clean source states, with code, environment, and
file hashes recorded in manifests. The reproduction guide documents
commands and historical provenance, including an earlier dirty source
state. Automated checks cover specified relationships; an unregistered
interpretation can still escape them.
